\section{Related Work}
\label{sec:related}

\para{LSH-based sparse attention.} \reformer \citep{kitaev2020reformer} replaces dense attention with random-projection LSH bucketing, achieving $O(L \log L)$ at the cost of forcing $Q = K$ and requiring training from scratch. \smyrf \citep{daras2020smyrf} relaxes the $Q=K$ constraint with an asymmetric transformation so distance in transformed space tracks inner product, enabling drop-in replacement of pretrained attention layers. \magicpig \citep{chen2024magicpig} argues that LSH-Top-K is a \emph{biased} estimator that collapses on aggregation tasks (\ruler-cwe/fwe), and instead uses LSH-sampling proportional to attention weight; it places hash tables on CPU to avoid GPU-unfriendly hash operations. \hashevict \citep{liu2024hashevict} uses pre-attention LSH eviction with binary codes and Hamming distance for the \kvcache. {\bf Unlike all of these, our index is keyed by token identity, not by query embedding}; it caches a token's historical attention behavior for reuse across occurrences, an axis prior LSH-attention work does not exploit. \hma also avoids the LSH-Top-K aggregation failure by using attention-weight thresholding (\eg storing all positions with weight $> \tau$) rather than strict top-$k$ selection.

\para{KV-cache eviction and selection.} \streaming \citep{xiao2023streamingllm} preserves a small set of attention sinks and a sliding window. \hto \citep{zhang2023h2o} ranks tokens by accumulated attention and evicts the rest. \quest \citep{tang2024quest} uses page-level min/max bounds to cheaply select Top-K pages, yielding $7\times$ attention speedup at long context. \cite{liu2024retrievalattention} identifies a query-key out-of-distribution problem when ANN indexes are built only over keys, and solves it with query-aware index construction. We compare \hma to \streaming, \hto, and a random-selection control at matched effective budget on \gpttwo-small. We do not reproduce \quest or \magicpig end-to-end at this scale; the comparison we draw is at the \emph{primitive} level (token-identity lookup vs.\ query-aware hashing), not in the systems regime.

\para{KV-cache quantization.} \kivi \citep{liu2024kivi} achieves 2-bit \kvcache compression with per-channel quantization for keys (because outlier channels persist in K) and per-token quantization for values (because V is mixed by attention). \kvquant \citep{hooper2024kvquant} extends to per-channel pre-RoPE outlier handling. \cite{kvtuner2025} learns layer-wise mixed-precision policies. {\bf None of these methods drive bit-width from observed attention frequency}; bit allocation is either static or set offline. Our \hmaquant policy directly reuses the lookup table's hit count to identify which positions deserve more bits, and we show this beats uniform 4-bit and random allocation at matched memory.

\para{Trainable sparse attention.} \nsa \citep{deepseek2025nsa} demonstrates that sparse attention can match or exceed full attention when natively pretrained, and explicitly notes that hard hash-based selection (as in \magicpig) is non-differentiable, blocking gradient flow. Our \hma in its current form is a post-hoc, inference-time primitive. A trainable variant would require relaxing the threshold gate, which we leave to future work.

\para{Other approximate attention.} Linear-attention methods \citep{performer,linear-transformers} replace softmax with separable kernels, fundamentally different in mechanism. Routing Transformer \citep{routing-transformer} clusters Q and K via $k$-means. \cite{ribar2023sparq} prunes Q channels for cheap approximate scores. \cite{loki2024} ranks tokens in a low-rank subspace of K. These methods all hash or summarize \emph{at query time}; \hma is novel in caching \emph{at prefill time, indexed by token identity}.
